\chapter{Dyskusja}

\section{Interpretacja uzyskanych wyników}
Analiza rezultatów przedstawionych w poprzednim rozdziale wyraźnie wskazuje, że związek pomiędzy danymi telemetrycznymi (takimi jak profil wygenerowanej mocy czy tętno w różnych strefach intensywności) a poziomem funkcjonalnego progu mocy (FTP) charakteryzuje się nieliniowością. Lepsza skuteczność modeli opartych na zespołach drzew decyzyjnych (m.in. Random Forest, XGBoost) w porównaniu do bazowej regresji liniowej wynika w głównej mierze ze zdolności estymatorów nieliniowych do przechwytywania złożonych interakcji pomiędzy zmiennymi. 

Szczególne znaczenie dla dokładności predykcji mają cechy agregujące informacje o pracy wykonanej na dłuższych dystansach i w dłuższych oknach czasowych. Zależność mocy od tętna (np. wskaźnik efektywności tętna) oraz badanie zjawiska dryfu tętna (ang. \textit{cardiac drift}) w trakcie jednostajnego wysiłku pozwalają uchwycić zmiany sprawności metabolicznej i układu sercowo-naczyniowego, które ściśle korelują z parametrem FTP.

\section{Znaczenie modeli uczenia maszynowego w estymacji FTP}
Zastosowanie algorytmów uczenia maszynowego zmienia tradycyjny paradygmat monitorowania dyspozycji kolarza. Klasyczne podejścia zakładają konieczność wykonywania okresowych, wysoce wyczerpujących testów sprawdzających (np. 20-minutowy test FTP w terenie) \cite{sitko2023_ac20to60_update}. Rozwiązania analityczne pokazują, że informacje niezbędne do uaktualniania statusu formy zawarte są także w codziennych, nierzadko sub-maksymalnych treningach \cite{phasescycling2023_ftp_variability}. Odpowiednia inżynieria cech wejściowych oraz optymalizacja modeli numerycznych umożliwia ekstrakcję tego sygnału i regularne estymowanie obiektywnego wskaźnika wydolności tlenowej bez zaburzania struktury planu treningowego.

\section{Ograniczenia przeprowadzonego badania}
Mimo obiecujących rezultatów zaproponowanego w pracy podejścia, należy zwrócić uwagę na istotne ograniczenia zrealizowanych eksperymentów:
\begin{itemize}
    \item \textbf{Zaszumienie danych telemetrycznych:} Dane gromadzone w terenie przy użyciu komercyjnych mierników mocy i pulsometrów obarczone są błędami pomiarowymi, przerwami w transmisji czy artefaktami.
    \item \textbf{Kwestia wartości referencyjnych (Ground Truth):} Wykorzystane wartości FTP były nierzadko wyznaczane heurystycznie lub z danych zewnętrznych systemów i same w sobie mogą nie być idealnym odzwierciedleniem fizjologicznej dyspozycji. 
    \item \textbf{Uproszczenie etykietowania:} Założenie okna propagacji etykiety (np. $\pm 14$ dni) dookoła zidentyfikowanego testu jest znaczącym uproszczeniem, nie uwzględniającym mikrozmian i wahań dyspozycji występujących z dnia na dzień.
    \item \textbf{Braki w dodatkowych atrybutach:} Model operuje w głównej mierze na danych o pracy fizycznej. Brak w nim integracji z metrykami pochodzącymi z życia codziennego, takimi jak jakość i długość snu, regeneracja, poziom stresu, zbilansowanie diety czy warunki atmosferyczne (temperatura, wiatr). 
    \item \textbf{Zdolność do uogólniania:} Wyniki uzyskane na zebranej próbie badawczej (obejmującej ostatecznie 126 kolarzy) mogą nie generalizować się w pełni na osoby na skrajnie różnym poziomie zaawansowania sportowego bez zastosowania kalibracji osobniczej.
\end{itemize}

\section{Możliwości praktycznego zastosowania}
Rozwiązania podobne do wytworzonego modelu predykcyjnego mogą z powodzeniem zostać zaimplementowane w nowoczesnych ekosystemach sportowych. Należą do nich przede wszystkim:
\begin{itemize}
    \item \textbf{Platformy analityczne i aplikacje treningowe:} Oprogramowanie może na bieżąco monitorować rozwój formy zawodnika na podstawie wgrywanych plików treningowych.
    \item \textbf{Automatyczna aktualizacja stref:} Adaptacyjna modyfikacja stref treningowych oparta na ciągłej, pasywnej ocenie formy, minimalizująca błędy w zadawaniu intensywności zadań.
    \item \textbf{Wsparcie dla trenerów:} Narzędzie wczesnego ostrzegania informujące o stagnacji bądź pogarszających się parametrach fizjologicznych na długo przed planowanym sprawdzianem wydolności.
\end{itemize}

\section{Kierunki dalszych badań}
Rozwój zaproponowanego systemu jest bardzo szeroki. Przyszłe badania mogłyby skoncentrować się na implementacji bardziej zaawansowanych struktur modelowych, takich jak sieci neuronowe dla danych sekwencyjnych (LSTM, GRU) czy architektury oparte na mechanizmie uwagi (Transformer), które pozwalają bezpośrednio analizować zjawiska w szeregach czasowych bez wymogu ręcznej ekstrakcji cech.

Kolejnym naturalnym krokiem jest wykorzystanie większych korpusów danych przestrzenno-czasowych z perspektywy kilku kompletnych sezonów startowych. Wskazana byłaby także rygorystyczna walidacja badawcza zestawiona z pomiarami w kontrolowanych warunkach laboratoryjnych (np. testami poziomu kwasu mlekowego) oraz prace nad algorytmami wysoce spersonalizowanymi, dostrojonymi bezpośrednio pod profil fizjologiczny pojedynczego zawodnika.
